\subsection{What Our Meetings Are For}

Installment 3: “Folkebladet,” March 08, 1916

There could be much to say about the reason why our meetings are less well attended than before; but I shall pass over that this time and look somewhat at our meetings as we generally have them.

At times the thought has taken hold of me that our meetings are not always what they ought to be. And on this account the question has arisen: what is the deficiency, or what are the deficiencies, in several of our meetings?

As will be understood, I proceed from the assumption that we intend something, that we will something, when we hold meetings.

For it is also conceivable that meetings may be held without anything real being intended by them.

But what is it, then, that we will?

The answer may perhaps be given: we wish, if possible, to derive some profit for faith and thought. We intend that the individual, through the word and the Spirit, may come nearer to God. Among other things, we intend through our meetings that those who gather may be influenced by the truth toward a deeper recognition of sin and guilt before God, and toward a more childlike and trusting appropriation of grace, of salvation.

And finally, we intend to impart Christian knowledge, deeper insight into the Lord’s word and the Lord’s plan of salvation for the individual and for the congregation, and a clearer understanding of our life and of the task we have been set to accomplish here in life. In a word: we would so gladly be made partakers of a greater fullness of God’s Spirit.

As may be seen, the goal here set as the result of a meeting is rather high. But not too high. It is not unattainable, even though it is not easy to reach. Yet the difficulty involved in attaining this high goal ought neither to frighten us away from holding such meetings nor, in faintheartedness, cause us to lower the demand or comfort ourselves with the cheap consolation that it matters little whether the goal is reached or not.

No, let the high goal stand before our thought and our longing in all its beauty and greatness; and let us, with all the means God himself has entrusted to us, labor and pray that it may be attained.

The truth is, namely, that it can be attained. For of all that is certain, nothing is more certain than this, that “the Lord holds his people dear.” No gift is the Lord more willing to give his people than the gift of the Spirit. And when his people gather as children around their Father’s table to be sanctified and made fit for their work for him, the Lord is willing to bless.

It is, however, a fact that the experience of many concerning the profit of such meetings is this: that it has often been so small. And by profit I do not mean primarily the more or less powerful stirring which, for many, is the only profit they desire. No, I think rather of what is called spirit in a meeting, of whether we succeed in leading one another deeper into the word, and finally of the abiding effects of the word spoken during the meeting.

What, then, is the reason that the profit is so often so small? What is the reason that many of our meetings are so ineffectual, so poor in spirit and in profit for the participants?

The causes are no doubt many, perhaps. And I do not presume to think that I possess the ability to point out all the hindering causes on account of which several of our meetings bear so little fruit. But if I could succeed in drawing the attention of our Free Church people to this matter, and if I could moreover point out some of these hindering causes and, if possible, do ever so little toward their removal, the purpose of these words would be attained.

There are certain things in connection with our larger meetings that I wish to mention, but I am uncertain which I should mention first. I therefore point to the hindrances I believe I see in the order in which they present themselves to my thought.

The first great deficiency in at least most of our meetings, I believe, is this: they are far too little prepared through prayer.

I know well that some who read this will object that this is not so, but that believers, prior to the holding of a meeting, pray for God’s blessing upon it. Yes, at almost all the meetings I have attended, I myself have heard speakers say that they have asked the Lord to “save, if only one soul, at the meeting.” And this, then, is supposed to be proof that the meeting has been well prepared through prayer? To such a claim I will say, quietly but firmly: no!

For in the first place, very few of those who pray for a meeting intend anything definite by their prayer. They pray because they believe it is their duty to do so. Or, if they do intend something by their prayer, there are many who do not know what they intend. And how, then, can God answer such a prayer?

There are no doubt those who will contradict me here also and answer that they both intend something and know what they intend. For, they say, we ask the Lord to save the unsaved, and we pray that God’s children may be edified; and then surely we must intend something and know what it is we intend.

Well, yes, I myself have many times prayed in the same manner, but I fear that, nevertheless, I have not truly been clear about what I was asking. What is meant when in prayer we say to God: you must save sinners at this meeting?

I assume that the very great majority who pray in this manner think chiefly of a spiritual stirring perceptible to everyone during the meeting, especially among those who were assumed until then to have been indifferent.

Yes, someone says, but is this not good?

Yes, of course it is. And more than once during the time in which I have taken part in meetings, I myself have experienced the great joy of witnessing such a stirring. And yet—it is not decisive.

Spiritual stirring is not synonymous with “being saved.” What is decisive—the fact that a soul in its inmost being turns from the world and sin to God—is not proved by an hour of emotional stirring during a meeting. The stirring that is perceived may be deep and true, but it may also be superficial and deceptive.

It is therefore not certain that all who pray that a “soul may be saved” during the meeting are thinking of, and truly know, what they are asking. They may pray earnestly and mean well, but they do not think deeply enough, not biblically enough. For as a rule, by the expression “the salvation of a soul,” one thinks only of being allowed, at the meeting, to witness one or more persons, with tears and perhaps with a request for intercession, show openly that they have been seized by the effect of the word. And this—as already said—is good and well so far as it reaches. But with many, with most, it does not reach far enough.

Salvation, or perhaps more properly, being made a partaker of salvation, is something else and something more—not merely religious moods, not merely anguish over having acted foolishly, and still less terror of hell, but the heart’s surrender to God, the will at rest in harmony with his will, so that the sinner in the inmost depth of his being says: “I will rise and go to my Father.”

By this I have not said that we should not desire and pray for spiritual stirring during a meeting. No, we can hardly have too much of it—at least not of that stirring which is brought about by a living and Spirit-filled proclamation of God’s word.

For it is a truth confirmed by experience that the word most often makes its way through feeling to the will.

There may, of course, be such natures whose life of the will is best influenced through thought, but they are exceptions, not the rule. Therefore let us rejoice that spiritual stirring is perceptible during a meeting, but let us not overvalue it.